\begin{figure}[htbp]
\centering
\begin{tikzpicture}
\begin{axis}[
    width=0.96\textwidth,
    height=4.5cm,
    grid=both,
    grid style={dotted, gray!40},
    xlabel={Messzeit $t$ [Sekunden]},
    ylabel={Anzahl Replikate},
    xmin=0, xmax=10,
    ymin=-0.2, ymax=2.5,
    ytick={0, 1, 2},
    legend pos=south east,
    legend style={font=\sffamily\scriptsize, fill=white, fill opacity=0.9},
    label style={font=\sffamily\footnotesize},
    tick label style={font=\sffamily\scriptsize},
    title style={font=\sffamily\bfseries\footnotesize},
]

% Soll-Zustand (spec.replicas)
\addplot[
    const plot,
    color=orange!85!black,
    mark=square*,
    mark size=1.3pt,
    line width=1.1pt,
    dashed
] table[x=elapsed_s, y=desired_replicas, col sep=comma] {daten/exp3-timeseries.csv};
\addlegendentry{Soll-Replikate (\texttt{spec.replicas})}

% Ist-Zustand (status.readyReplicas)
\addplot[
    const plot,
    color=blue!70!black,
    mark=*,
    mark size=1.3pt,
    line width=1.1pt
] table[x=elapsed_s, y=ready_replicas, col sep=comma] {daten/exp3-timeseries.csv};
\addlegendentry{Bereite Replikate (\texttt{readyReplicas})}

\end{axis}
\end{tikzpicture}
\caption{Zeitlicher Verlauf der Soll- und Ist-Replikate bei manueller Skalierung auf 0 und automatischer Wiederherstellung durch Argo~CD Self-Heal (Durchlauf 1), eigene Messung.}
\label{fig:exp3-selfheal-verlauf}
\end{figure}
